\textbf{From recorded robot interaction to a forecasting test.}
(A) Camera 1 supplies training, validation and primary testing; camera 2 is used in a separate held-out viewpoint-transfer test. The wrist camera is retained for inspection. Camera placement is schematic, with one camera stream per model run.
(B) Three recorded observations at native frames 0, 5 and 10, two past command blocks and five supplied future blocks condition an offline forecast. Each block contains five seven-dimensional recorded commands. Five successive future feature vectors are predicted; the drawing shows the fifth, at frame 35.
(C) The corresponding recorded future image enters only the evaluator. Both paths use the same frozen DINOv2-small encoder, $2\times2$ pooling and training-camera-1 feature statistics. Predictions are compared with the fixed 1,536-dimensional targets using feature MSE and cosine error. The section reports all registered horizons and both camera populations separately.
Photos show the first eligible test episode by lexicographic ID, selected independently of outcomes. All views retain their full frame and publisher-provided blurring. The camera illustration and feature tiles are schematic, while photographs are unchanged DROID observations (CC BY 4.0). This is the historical context-model study; the current spatial architecture is in Figure~\ref{fig:editorial-spatial-method}.
